\documentclass{article}
\usepackage{amsmath, amssymb, amsthm}
\usepackage{hyperref}
\usepackage{geometry}
\geometry{margin=1in}

\title{Erd\H{o}s Problem \#1017}
\date{Last updated: 2025-09-12}
\author{Source: erdosproblems.com}

\begin{document}
\maketitle

\noindent\textbf{Status:} OPEN \quad \textbf{No prize}

\noindent\textbf{Tags:} graph theory

\medskip
\noindent\textbf{Problem Statement:}

Let $f(n,k)$ be such that every graph on $n$ vertices and $k$ edges can be partitioned into at most $f(n,k)$ edge-disjoint complete graphs. Estimate $f(n,k)$ for $k>n^2/4$.

\bigskip
\noindent\textbf{Remarks:}

The function $f(n,k)$ is sometimes called the clique partition number.

Erd\H{o}s, Goodman, and P\'{o}sa \cite{EGP66} proved that $f(n,k)\leq n^2/4$ for all $k$ (and in fact the complete graphs can be taken to be edges and triangles), which is best possible in general, as witnessed for example by a complete bipartite graph. In \cite{Er71} Erd\H{o} asks vaguely whether this result can be 'sharpened' for $k>n^2/4$.

Lov\'{a}sz \cite{Lo68} proved that every graph on $n$ vertices and $k$ edges is the union of $\binom{n}{2}-k+t$ complete graphs, where $t$ is maximal such that $t^2-t\leq \binom{n}{2}-k$, but without the assumption that the complete graphs are edge disjoint. Lov\'{a}sz's result is sharp in many cases.

If $k>n^2/4$ and the graph contains no $K_4$ then this is equivalent to finding the minimum number of edge disjoint triangles. This special case was also asked about by Erd\H{o}s. A complete answer was provided by Gy\"{o}ri and Keszegh \cite{GyKe17}, who proved that every $K_4$-free graph with $n$ vertices and $\lfloor n^2/4\rfloor+m$ edges contins $m$ pairwise edge disjoint triangles.

See also [184] for an analogous problem decomposing into edges and cycles, and [583] for decomposing into paths. The clique partition problem for chordal graphs is the subject of [81].

\bigskip
\noindent\textbf{References:}

[EGP66] Erd\H{o}s, Paul and Goodman, A. W. and P\'{o}sa, Lajos, The representation of a graph by set intersections. Canadian J. Math. (1966), 106-112.
 
 [Er71] Erd\H{o}s, P., Some unsolved problems in graph theory and combinatorial analysis. Combinatorial Mathematics and its Applications (Proc.
Conf., Oxford, 1969) (1971), 97-109.
 
 [GyKe17] Gy\H{o}ri, Ervin and Keszegh, Bal\'{a}zs, On the number of edge-disjoint triangles in {$K_4$}-free
graphs. Combinatorica (2017), 1113--1124.
 
 [Lo68] Lov\'{a}sz, L., On covering of graphs. Theory of Graphs (Proc. Colloq., Tihany, 1966) (1968), 231-236.

\bigskip
\noindent\small{Source: \url{https://www.erdosproblems.com/1017}}
\end{document}
